% Preámbulo común del material docente · XeLaTeX
\documentclass[11pt,a4paper]{article}
\usepackage[margin=22mm,top=20mm,bottom=22mm]{geometry}
\usepackage{fontspec}
\setmainfont{Avenir Next}[BoldFont={Avenir Next Demi Bold}, BoldItalicFont={Avenir Next Demi Bold Italic}]
\setsansfont{Avenir Next}
\setmonofont{Menlo}[Scale=0.88]
\newfontfamily\symfont{Apple Symbols}
\newfontfamily\unifont{Arial Unicode MS}
\usepackage{unicode-math}
\setmathfont{latinmodern-math.otf}
\usepackage{polyglossia}\setdefaultlanguage{spanish}
\usepackage{xcolor,booktabs,tabularx,enumitem,parskip,microtype,titlesec,fancyhdr}
\usepackage[most]{tcolorbox}
\usepackage[hidelinks]{hyperref}
\emergencystretch=2em

\definecolor{tinta}{HTML}{1F2933}
\definecolor{acento}{HTML}{B7410E}
\definecolor{suave}{HTML}{F4F1EC}
\definecolor{gris}{HTML}{6B7280}
\definecolor{linea}{HTML}{D6D0C6}
\color{tinta}

\titleformat{\section}{\Large\bfseries\color{acento}}{\thesection}{0.6em}{}
\titleformat{\subsection}{\large\bfseries}{\thesubsection}{0.6em}{}
\titlespacing*{\section}{0pt}{1.6em}{0.5em}
\titlespacing*{\subsection}{0pt}{1.1em}{0.3em}
\setlist{nosep, leftmargin=1.4em, itemsep=2pt}
\renewcommand{\arraystretch}{1.25}

% bloque de órdenes de terminal
\newtcblisting{orden}{listing only, colback=suave, colframe=linea, boxrule=0.4pt, arc=2pt,
  left=6pt, right=6pt, top=4pt, bottom=4pt, listing options={basicstyle=\ttfamily\small, breaklines=true, columns=fullflexible, keepspaces=true}}
% nota destacada
\newtcolorbox{nota}[1][Nota]{colback=white, colframe=acento, boxrule=0.6pt, arc=2pt, title=#1,
  fonttitle=\bfseries\small, coltitle=white, colbacktitle=acento, left=6pt, right=6pt}
% preguntas
\newcommand{\pregunta}[2]{\item[\textbf{(#1)}] #2}
\newenvironment{preguntas}{\begin{description}[leftmargin=2.4em, style=nextline, itemsep=4pt]}{\end{description}}
% caja para que el estudiante escriba
\newtcolorbox{respuesta}[2][2.2cm]{colback=white, colframe=linea, boxrule=0.5pt, arc=1.5pt, height=#1,
  title={#2}, fonttitle=\bfseries\small, coltitle=acento, colbacktitle=suave, titlerule=0pt,
  toptitle=1.5pt, bottomtitle=1.5pt, left=6pt, right=6pt, top=3pt, bottom=3pt, before skip=4pt, after skip=8pt}
\newcommand{\ok}{{\symfont ✔}}
\newcommand{\nok}{{\symfont ✘}}
\newcommand{\qed}{{\symfont ∎}}
\newcommand{\codigo}[1]{\texttt{#1}}

\pagestyle{fancy}\fancyhf{}
\renewcommand{\headrulewidth}{0pt}
\fancyfoot[C]{\small\color{gris}\docpie\ · página \thepage}

\newcommand{\cabecera}[3]{%
  {\Huge\bfseries\color{acento}#1\par}\vspace{4pt}
  {\large #2\par}\vspace{2pt}
  {\color{gris}#3\par}\vspace{8pt}
  {\color{linea}\hrule height 0.6pt}\vspace{10pt}}

\newcommand{\docpie}{Laboratorio 3 · notas del docente}
\begin{document}
\cabecera{Laboratorio 3 · notas del docente}{Interpolación con un agente de RL · 90 minutos}{Resultados medidos el 6 de octubre de 2026 (Linux, Python 3.12, semilla 1 salvo que se diga otra cosa). Todo se reproduce con \codigo{docs/curso/lab03/experimentos.py} y \codigo{graficar.py}. Usa el proyecto 07; en el arco estaba como lab 5 y se adelantó porque la clase del 1 de octubre cubrió Taylor y Lagrange. El laboratorio de Taylor (proyecto 03) sigue pendiente.}

\section{Plan de la sesión}
\begin{tabularx}{\linewidth}{@{}l X l@{}}
\toprule
\textbf{min} & \textbf{bloque} & \textbf{modo} \\
\midrule
0–10 & descargar el ZIP de hoy, instalar 07 con \codigo{matplotlib}, \codigo{--no-tui} (20 s) & guiado \\
10–25 & basales a mano (a)–(d), sin IA; revisar con \codigo{basales} & parejas, papel \\
25–40 & \codigo{graficar.py}: ajustar \codigo{PUNTOS}, \codigo{RANGO}, \codigo{NODOS}; PIB (e)–(g) & parejas \\
40–55 & agente con animación, dos semillas, informe (h)–(i) & parejas \\
55–70 & \codigo{familias}, \codigo{solo\_equi}, \codigo{sin\_runge}, \codigo{recompensa} (j)–(l) & parejas + plenaria corta \\
70–85 & \codigo{dibujar}, \codigo{grafo} (m)–(n); \codigo{mi\_grafo} (o) & parejas + plenaria \\
85–90 & cierre y entrega & expositivo \\
\bottomrule
\end{tabularx}

Las ideas de la sesión son tres: la suma de valores absolutos de los basales es lo que amplifica los errores (de (c) a (g) y a (h)); lo que el agente aprende depende de los problemas que ve ((k)); y el verificador solo sabe lo que dicen las flechas ((n) y (o)). Si se va el tiempo, cortar (i), (k) y la tarea opcional, en ese orden. No cortar (c), (g) ni \codigo{mi\_grafo}.

\section{Tiempos medidos}
\begin{tabularx}{\linewidth}{@{}X l@{}}
\toprule
\textbf{orden} & \textbf{tiempo} \\
\midrule
\codigo{python -m interprl --no-tui --seed 1} & 17 s (fase 1: 17 s; fase 2: 0,4 s) \\
\codigo{python -m interprl --seed 1} (animación) & $\sim$2 min \\
\codigo{basales}, \codigo{dibujar}, \codigo{grafo}, \codigo{mi\_grafo}, \codigo{graficar.py} & $<$ 2 s \\
\codigo{familias} & 2 s \\
\codigo{solo\_equi}, \codigo{sin\_runge} & 15–16 s (dos entrenamientos de 2000 gen.) \\
\codigo{recompensa} & 23 s (tres entrenamientos) \\
\bottomrule
\end{tabularx}

En los PCs del laboratorio puede ser el doble. Avisar que \codigo{--no-tui} no está colgado.

\section{Respuestas y cifras}
\begin{preguntas}
\pregunta{a}{$L_0 = t(t-1)/2$, $L_1 = 1-t^2$, $L_2 = t(t+1)/2$.}
\pregunta{b}{$L_0 = (t-1)(t-3)/3$, $L_1 = -t(t-3)/2$, $L_2 = t(t-1)/6$. El error típico es el signo de $L_1$: el denominador es $(1-0)(1-3)=-2$.}
\pregunta{c}{Lo que imprime \codigo{basales}:}
\end{preguntas}
\begin{center}
\begin{tabular}{ccccccc}
\toprule
\textbf{nodos} & $t$ & $L_0$ & $L_1$ & $L_2$ & \textbf{suma} & \textbf{suma de $|L_k|$} \\
\midrule
$-1,0,1$ & $-1/2$ & $3/8$ & $3/4$ & $-1/8$ & 1 & $5/4$ \\
$-1,0,1$ & $1/2$ & $-1/8$ & $3/4$ & $3/8$ & 1 & $5/4$ \\
$0,1,3$ & $2$ & $-1/3$ & $1$ & $1/3$ & 1 & $5/3$ \\
$-1,0,1$ & $2$ & $1$ & $-3$ & $3$ & 1 & $7$ \\
\bottomrule
\end{tabular}
\end{center}
La constante de Lebesgue de $-1,0,1$ en $[-1,1]$ es exactamente $5/4$, alcanzada en $t=\pm 1/2$: el número que calcularon a mano. Vale la pena decirlo en voz alta.

\begin{preguntas}
\pregunta{d}{Una versión aceptable, que es también la solución de \codigo{mi\_grafo} (7 pasos; el agente la ordena en la generación 5):}
\end{preguntas}
\begin{center}
\begin{tabularx}{\linewidth}{@{}l X l@{}}
\toprule
\textbf{clave} & \textbf{afirmación} & \textbf{deps} \\
\midrule
H1 & los nodos son distintos & — \\
DEF & $L_k(t)=\prod_{j\neq k}(t-x_j)/(x_k-x_j)$, grado $n$ & H1 \\
DELTA & $L_k(x_j)=1$ si $j=k$, 0 si no & DEF \\
S & $S(t)=\sum_k L_k(t)$ tiene grado $\le n$ & DEF \\
SNODOS & $S(x_j)=1$ para todo $j$ & S, DELTA \\
UNIQ & un polinomio de grado $\le n$ con $n+1$ ceros es nulo & H1 \\
QED & $S-1$ tiene grado $\le n$ y $n+1$ ceros, luego $S\equiv 1$ & S, SNODOS, UNIQ \\
\bottomrule
\end{tabularx}
\end{center}
Distractor natural: «$S$ vale 1 en los $n+1$ nodos, luego $S\equiv 1$» (le falta el grado $\le n$; la flecha de S a QED es la que lo descarta).

\begin{preguntas}
\pregunta{e}{Vietnam: $P(2013)=216{,}3$ frente a 213,7 real ($+1{,}2\,\%$); $P(2015)=247{,}1$ frente a 239,3 ($+3{,}3\,\%$). La serie es suave entre los nodos y los basales en $\pm1/2$ están entre $-1/8$ y $3/4$: no amplifican. 2015 es peor porque el crecimiento en dólares casi se detiene ese año ($233{,}5\to 239{,}3$).}
\pregunta{f}{Colombia: $P(2020)=317{,}8$ frente a 270,3 real ($+47{,}5$; $+17{,}6\,\%$). El polinomio solo sabe lo que dicen los nodos, y la caída de 2020 no dejó rastro en 2018, 2019 ni 2021. La fórmula del error necesita $f'''$ acotada y conocida; con datos no hay $f$, y un choque de un año hace que cualquier $f$ razonable tenga $f'''$ enorme. La fórmula dice cuándo confiar, no cuánto se equivoca uno con datos reales.}
\pregunta{g}{China: $P(2022)=20\,179{,}9$ frente a 17\,881,8 ($+12{,}9\,\%$); $P(2023)=22\,932{,}0$ frente a 17\,794,8 ($+28{,}9\,\%$). Basales en $t=3/2$: $3/8$, $-5/4$, $15/8$ (suma de absolutos 3,5); en $t=2$: $1$, $-3$, $3$ (suma 7). Un 1\,\% del dato de 2021 (178,2) se multiplica por $L_2(2)=3$: la estimación de 2023 se mueve 534,6. Interpolar no es pronosticar.}
\end{preguntas}

\textbf{Hitos} (\codigo{experimentos.py hitos --seed N}; idénticos a los de la animación con la misma semilla):
\begin{center}
\begin{tabular}{lcc}
\toprule
\textbf{hito} & \textbf{seed 1} & \textbf{seed 2} \\
\midrule
primera aproximación correcta & 80 & 12 \\
descubrió el criterio de parada & 80 & 111 \\
su regla coincide con la teoría (8 estados, estable 100 gen.) & 526 & 542 \\
prefiere $\gamma=1$ de forma estable & 726 & 572 \\
primera demostración completa & 186 & 215 \\
primera demostración mínima & 463 & 368 \\
\bottomrule
\end{tabular}
\end{center}
Con seed 2 aparece además «$\ge 95\,\%$ de aciertos» en la 2490; con seed 1 no. Al final: seed 1, 89\,\% de éxito en las últimas 100, 2187 correctas, 792 prematuras, 21 agotadas; seed 2, 80\,\%, 2186, 782, 32.

\begin{preguntas}
\pregunta{h}{Seed 1: $Q(\gamma=1)=-2{,}8$ (2584 usos) frente a $Q(\gamma=0)=-13{,}4$ (110 usos). $Q(\gamma)$ es el retorno medio de una partida con esa familia; con equiespaciados las funciones tipo Runge acaban en entrega prematura o agotadas. Tabla de $\Lambda$ del informe: con $n=16$, equiespaciados 934,5, Chebyshev 2,7. Es lo mismo que la última fila de (c): basales grandes amplifican, y con equiespaciados crecen como $2^n$.}
\pregunta{i}{Cuando los interpolantes convergen geométricamente, $|p_n-p_{n-1}|$ es del orden del error de $p_{n-1}$ y acota bien el de $p_n$. Engaña cuando todavía no hay convergencia o el error se estanca. Costo: 792 entregas prematuras de 3000 con seed 1. Es el problema de Colombia: lo que no se ve en las muestras no existe para el método.}
\pregunta{j}{\codigo{familias} (120 problemas, máximo 40 nodos): Runge $1/(1+ax^2)$, equiespaciados 0/31, Chebyshev 31/31 con 29,2 nodos; polo $1/(x-c)$, 27/29 (22,5) frente a 29/29 (16,5); raíz, 17/17 (14,9) frente a 17/17 (11,2); enteras, 43/43 (10,4) frente a 43/43 (9,5). \codigo{solo\_equi} (2000 gen.): éxito 85\,\% $\to$ 25\,\%, prematuras 653 $\to$ 1505, agotadas 14 $\to$ 7. Con Runge cada nodo equiespaciado sube el error: el progreso es negativo en cada paso y seguir cuesta más que el $-10$ de entregar mal. Rendirse es lo racional con esos nodos.}
\pregunta{k}{\codigo{sin\_runge}: éxito 94\,\%, $Q(\gamma)$ entre $+0{,}7$ y $+7{,}5$, preferida $\gamma=0{,}75$. Ya no hay una familia que gane con claridad (las diferencias están dentro del ruido del bandido, $\alpha=0{,}05$). Lo aprendido en la parte 4 no es «Chebyshev es mejor» como teorema, sino «con estos problemas, Chebyshev da más puntos». El teorema es sobre los nodos; la preferencia del agente aparece solo cuando los problemas la hacen importar. La distribución de problemas es parte de la especificación, como la recompensa.}
\pregunta{l}{Original: 85\,\%, 1333/653/14, $\gamma=1$. «Entregar mal cuesta 0»: 0\,\%, 1 correcta y 1999 prematuras: entrega siempre en el primer paso. Una partida que trabaja hace unos 18 nodos ($-18$), el progreso telescopa a $\log_{10}(\text{error inicial}/\text{error final})\approx 5$ y gana $+10$: unos $-3$. Entregar ya da 0. Es el (f) del lab 1. «Cada nodo gratis»: 87\,\%, $\gamma=1$, casi nada cambia: el descuento 0,95 hace que el $+10$ valga más cuanto antes y, pasado el error de máquina, añadir nodos no da progreso. También probado \codigo{--entregar-mal -1 --nodo 0}: 85\,\%, $\gamma=1$.}
\pregunta{m}{Si los $n+1$ nodos se juntan en $x_0$, las condiciones pasan a ser $f(x_0), f'(x_0),\dots,f^{(n)}(x_0)$ (Hermite en el límite), el interpolante es el polinomio de Taylor y $w(x)\to(x-x_0)^{n+1}$. Taylor es el caso en que toda la información está en un punto: la idea con la que se abrió la clase del 1 de octubre.}
\pregunta{n}{Si $x=x_i$, $w(x)=0$ y $g$ ni siquiera está definida; FIX trata ese caso aparte (ambos lados son 0). Sin la flecha FIX $\to$ AUX, la demostración de 12 pasos que acepta el verificador divide por cero sin decirlo. Responsable: quien escribió las \codigo{deps}.}
\pregunta{o}{Libre. Distractores buenos: el del grado (arriba); «cada $L_k$ está entre 0 y 1, luego la suma es $\le n+1$» (falso: la fila $t=2$ de (c)); «$S=1$ porque los basales son una partición de la unidad» (circular).}
\end{preguntas}

\textbf{Párrafo final.} Que mencionen: rellenar datos entre muestras de algo suave (sirve); no sirve para choques que no tocan los nodos (Colombia) ni para extrapolar (China); más nodos equiespaciados no arreglan nada (Runge). Bonus si conectan con $\Lambda$.

\textbf{Tarea opcional.} Equiespaciados con 5 nodos: $\Lambda=2{,}2078$. Interiores en $\pm0{,}7$: $\Lambda=1{,}7727$. \codigo{search --n 4} (1 s) encuentra $\pm0{,}620911$ con $\Lambda=1{,}55949$, por debajo del Chebyshev extendido (1,57017).

\section{Trampas}
\begin{itemize}
\item \textbf{El ZIP tiene que ser de hoy.} Si lo descargaron antes, no trae \codigo{docs/curso/lab03}.
\item \codigo{matplotlib} va en la misma línea de instalación. Si falta, \codigo{graficar.py} dice cómo instalarlo.
\item En \codigo{basales} los nodos van con \codigo{=}: \codigo{--nodos=-1,0,1}. Sin \codigo{=}, el \codigo{-1} parece una opción.
\item \codigo{graficar.py} se edita. Si alguien lo rompe, se recupera del ZIP o con \codigo{git checkout}.
\item Con \codigo{PUNTOS = [..., 2]} y \codigo{RANGO = None}, la línea en $t=2$ existe pero cae fuera del eje: es intencional, es el ajuste 2.
\item El PIB está en dólares corrientes: mezcla producción, inflación y tasa de cambio. La caída de Colombia entre 2014 y 2015 (381,2 $\to$ 293,5) es sobre todo devaluación del peso.
\item \codigo{pib\_tres\_paises.csv} está congelado (Banco Mundial vía \codigo{datasets/gdp}, CC BY 4.0, hasta 2023). \codigo{graficar.py --reconstruir} lo regenera.
\item \codigo{mi\_grafo.py} sin tocar da un error a propósito: \qed{} sin deps. Es el primer mensaje que ven.
\end{itemize}
\end{document}
